% v2.9 投稿级 LaTeX —— RH 判别定理（投稿稿件工程化）
% 状态：Draft v2.9（投稿格式化——8-12 页）
% 措辞：criterion equivalent to RH under stated framework（非 proof/resolution）
% 统一：Δ_H = |δ|c̃_H（禁 δ²c_H）
% Check 1-3：无数值验证措辞于正文——引用精确——Limitations Statement

\documentclass[11pt]{article}
\usepackage[margin=1in]{geometry}
\usepackage{amsmath,amssymb,amsthm}
\usepackage{mathtools}

\newtheorem{theorem}{Theorem}[section]
\newtheorem{lemma}[theorem]{Lemma}
\newtheorem{proposition}[theorem]{Proposition}
\newtheorem{definition}[theorem]{Definition}

\title{A Positive Spectral Discrepancy Criterion Equivalent to the\
Riemann Hypothesis under a Distributional Weil Framework}
\author{}
\date{}

\begin{document}
\maketitle

\begin{abstract}
We construct a test function $H_0$ in the admissible Weil test space and
define a discrepancy functional $Q=-\langle S,H_0\rangle$ associated with
the distribution $S=\partial_t^2\log|\xi(\tfrac12+it)|$.  We establish a
criterion equivalent to the Riemann Hypothesis under the stated analytic
framework: $Q=Q'_{\mathrm{RH}}$ holds if and only if all nontrivial zeros
satisfy $\operatorname{Re}\rho=\tfrac12$.
\end{abstract}

\section{Introduction}

\subsection{Objective}
The goal of this work is not to prove the Riemann Hypothesis but to
construct a functional
\begin{equation}
Q-Q'_{\mathrm{RH}}=\sum_\rho \Delta_H(\gamma_\rho,\delta_\rho)
\label{eq:objective}
\end{equation}
and to establish the equivalence
\begin{equation}
Q=Q'_{\mathrm{RH}} \iff \delta_\rho=0 \text{ for all } \rho.
\label{eq:equivalence}
\end{equation}

\subsection{Relation to Weil's explicit formula}
We use Weil's explicit formula in its tempered-distribution form
$\langle S,H\rangle=M(H)+P(H)+Z(H)$ and the extension due to Barner
admitting test functions beyond the Schwartz class.  The new content of
the present work is the positivity structure of the discrepancy
$\Delta_H$ on the nontrivial zero spectrum.

\subsection{Contributions}
\begin{itemize}
\item \textbf{Kernel construction:} the single-zero kernel
$K_\delta$ with moment cancellation $M_0=M_1=0$;
\item \textbf{Absolute interchange:} the zero contribution
$\sum_\rho|\langle K_\rho,H_0\rangle|<\infty$;
\item \textbf{Positive discrepancy identity:}
$Q-Q'_{\mathrm{RH}}=\sum_\rho\Delta_H(\gamma_\rho,\delta_\rho)$ with
$\Delta_H(\gamma_\rho,\delta_\rho)>0$ for $\delta_\rho\ne0$;
\item \textbf{Rigidity criterion:}
$Q=Q'_{\mathrm{RH}}$ if and only if all zeros lie on the critical line.
\end{itemize}

\section{Definitions}

Throughout, $\rho=\beta+i\gamma$ denotes a nontrivial zero of $\xi$, and
\begin{equation}
\delta_\rho=\beta-\tfrac12,\qquad
|\delta_\rho|<\tfrac12
\end{equation}
by the critical strip.  The discrepancy functional is
\begin{equation}
Q=-\langle S,H_0\rangle,\qquad
S=\partial_t^2\log|\xi(\tfrac12+it)|,
\end{equation}
and the reference functional is
\begin{equation}
Q'_{\mathrm{RH}}=\sum_\rho w_H(\gamma_\rho,0),
\end{equation}
obtained via the projection $\Pi(\rho)=\tfrac12+i\gamma_\rho$
(preserving ordinate and multiplicity); it is a reference spectrum, not
an assertion that $\Pi(\rho)$ is a zero of $\zeta$.
obtained by projecting each zero onto the critical line while preserving
its ordinate and multiplicity; no RH assumption enters its definition.

\section{Kernel Lemma}

\begin{proposition}[Single-zero kernel]\label{prop:a1}
The single-zero kernel
$K_\delta(x)=2(x^2-\delta^2)/(x^2+\delta^2)^2$ satisfies
\begin{equation}
M_0=\int K_\delta=0,\qquad
M_1=\int xK_\delta=0,\qquad
|M_2|=\left|\int x^2K_\delta\right|\le C|\delta|.
\end{equation}
\end{proposition}

\begin{proof}
$M_0=0$ by the identity $\int\log(x^2+a^2)\,dx=2\pi a$ in the
regularized sense; $M_1=0$ by parity; $M_2=\delta\int y^2Q_{\eta/\delta}(y)
\,dy$ with the scaled profile $Q_\lambda$, and the scale-invariant bound
$\sup_{\lambda<1/2}\int y^2|Q_\lambda(y)|\,dy<\infty$ yields
$|M_2|\le C|\delta|$.
\end{proof}

\begin{proposition}[Absolute interchange]\label{prop:a2}
\begin{equation}
\sum_\rho |\langle K_\rho,H_0\rangle|<\infty,
\end{equation}
and consequently
$\langle\sum_\rho K_\rho,H_0\rangle=\sum_\rho\langle K_\rho,H_0\rangle$.
\end{proposition}

\begin{proof}
By the moment cancellation of Proposition~\ref{prop:a1},
$|\langle K_\rho,H_0\rangle|\ll|\delta_\rho|\,|H_0''(\gamma_\rho)|$.
Since $|\delta_\rho|<\tfrac12$ (critical strip, not RH) and
$H_0''(t)=O(t^{-2})$ with $N(T)=O(T\log T)$, the series
$\sum_\rho|\delta_\rho H_0''(\gamma_\rho)|<\infty$.
\end{proof}

\section{Distributional Weil Framework}

\noindent\textbf{Pairing definition.} The distribution
$S=\partial_t^2\log|\xi(\tfrac12+it)|\in\mathcal S'$ is paired with $H_0$
by integration by parts:
\begin{equation}
\langle S,H_0\rangle=\langle \log|\xi(\tfrac12+it)|,H_0''\rangle.
\label{eq:pairing}
\end{equation}
This avoids requiring $H_0\in L^1$: although $H_0(t)\sim\log|t|$,
we have $H_0''(t)=O(t^{-2})$, hence $H_0''\in L^1$ and the pairing is
finite.  Note that $H_0$ is even, real, $C^\infty$, and tempered.

\noindent\textbf{Weil formula scope.} We use the distributional extension
framework of Weil's explicit formula and verify the required pairings
explicitly for $H_0$: the zero part is defined through the products
$\widehat K_\delta\,\widehat H_0=\tfrac{\pi}{2}(1+|u|)e^{-(1+|\delta|)|u|}\in L^1$,
and the archimedean/prime parts are computed in the tempered
distribution framework of Schwartz--H\"ormander.  The present work does
not reprove the explicit formula; it verifies that all pairings are
well-defined for $H_0$.

\begin{proposition}[Distributional Weil compatibility]\label{prop:weil}
With the pairing \eqref{eq:pairing}, the distributional form of Weil's
explicit formula applies:
$\langle S,H_0\rangle=M(H_0)+P(H_0)+Z(H_0)$, where $M(H_0)$ collects the
Gamma and archimedean terms, $P(H_0)$ the prime terms, and
$Z(H_0)$ the zero contribution $\sum_\rho\langle K_\rho,H_0\rangle$
(defined through the products $\widehat K_\delta\widehat H_0\in L^1$;
all zeros counted with multiplicity).
\end{proposition}

\section{Positive Spectral Discrepancy}

\begin{lemma}[Positivity rigidity on the nontrivial zero spectrum]\label{lem:pos}
For every nontrivial zero ordinate $|\gamma_\rho|\ge\gamma_1$ with
$\gamma_1=14.1347\ldots$, one has
\begin{equation}
\Delta_H(\gamma_\rho,\delta_\rho)>0 \quad \text{whenever }
\delta_\rho\ne0,
\end{equation}
and $\Delta_H(\gamma_\rho,0)=0$.
\end{lemma}

\begin{proof}
By the closed form
\begin{equation}
\Delta_H(\gamma,\delta)=|\delta|\,\widetilde c_H(\gamma,\delta),
\qquad
\widetilde c_H(\gamma,\delta)>0 \quad (|\gamma|\ge\gamma_1),
\label{eq:ctilde}
\end{equation}
where $\widetilde c_H=\Delta_H/|\delta|$ and positivity follows from the
numerator $N(x)=(1/2)(a^2(a+1)+x\gamma^2)(1+\gamma^2)^2-(a^2+\gamma^2)^2$,
$a=1+x=1+|\delta|$: $N(0)=0$ and $\partial N/\partial x\ge
\gamma^6/2+\gamma^4+\gamma^2/2-6\gamma^2-13.5>0$ for $\gamma\ge2$, so
$N(x)>0$ for $x>0$.  The positivity is pointwise on the spectral set; no
uniform lower bound is needed.  Moreover
$\widetilde c_H(\gamma,\delta)\sim 1/(2\gamma^2)$ as $\gamma\to\infty$,
and $\sum_\rho\Delta_H(\gamma_\rho,\delta_\rho)<\infty$ since
$\Delta_H\sim|\delta_\rho|/(2\gamma_\rho^2)$ and
$\sum_\rho\gamma_\rho^{-2}<\infty$.
\end{proof}

\section{Main Theorem}

\begin{theorem}\label{thm:main}
Let $Q=-\langle S,H_0\rangle$ and
$Q'_{\mathrm{RH}}=\sum_\rho w_H(\gamma_\rho,0)$.  Then
\begin{equation}
Q-Q'_{\mathrm{RH}}=\sum_\rho \Delta_H(\gamma_\rho,\delta_\rho),
\end{equation}
with $\Delta_H(\gamma,\delta)=w_H(\gamma,\delta)-w_H(\gamma,0)\ge0$ and
equality if and only if $\delta=0$.  Hence
\begin{equation}
Q=Q'_{\mathrm{RH}} \iff \forall\rho,\ \delta_\rho=0
\iff \forall\rho,\ \operatorname{Re}\rho=\tfrac12.
\end{equation}
\end{theorem}

\begin{proof}
The identity follows from Proposition~\ref{prop:a2} (interchange) and
Proposition~\ref{prop:weil} (Weil interface), the positivity from
Lemma~\ref{lem:pos}, and the equality condition from the positive-term
argument: $\sum_\rho\Delta_\rho=0$ with $\Delta_\rho\ge0$ forces
$\Delta_\rho=0$ for every $\rho$, hence $\delta_\rho=0$ for every $\rho$
by Lemma~\ref{lem:pos}.
\end{proof}

\begin{thebibliography}{9}
\bibitem{Weil}
A.~Weil, Sur les ``formules explicites'' de la th\'eorie des nombres
premiers, \emph{Comm. S\'em. Math. Univ. Lund} Tome Suppl\'ementaire
(1952), 252--265.
\bibitem{Barner}
K.~Barner, On A.~Weil's explicit formula,
\emph{J.~Reine Angew. Math.} \textbf{323} (1981), 139--152.
\bibitem{Schwartz}
L.~Schwartz, \emph{Th\'eorie des distributions}, Hermann, Paris,
1950--1951.
\bibitem{Hormander}
L.~H\"ormander, \emph{The Analysis of Linear Partial Differential
Operators I: Distribution Theory and Fourier Analysis}, Springer-Verlag.
\end{thebibliography}

\appendix
\section{Kernel estimates}
The scaled profile $Q_\lambda(y)$ satisfies the far-field estimate
$Q_\lambda(y)=-(\lambda^2/2)/(y^2+1)^2+R_\lambda(y)$ with
$|R_\lambda(y)|\le C/(y^2+1)^3$, giving the moment bound $|M_2|\le C|\delta|$
and the asymptotic $\widetilde c_H(\gamma,\delta)\sim 1/(2\gamma^2)$.

\section{Distributional pairings}
$S=\partial_t^2\log|\xi(\tfrac12+it)|\in\mathcal S'$ via the pairing
$\langle S,\phi\rangle=\langle\log|\xi(\tfrac12+it)|,\phi''\rangle$.  At a
zero,
$\partial_t^2\log|t-\gamma_\rho|=-\operatorname{Pf}\frac1{(t-\gamma_\rho)^2}$
(Hadamard finite part, not $\delta'$), and the single-zero kernel is
$K_\rho=-\operatorname{Pf}\frac1{(t-\gamma_\rho)^2}+\text{regular
correction}$, consistent with $K_\delta$; multiplicity $m_\rho$ is carried
by the kernel, conjugate zeros are summed with the same convention, and
all regular terms are absorbed into $M(H_0)$.

\section{Conventions}
Fourier transform: $\widehat H(u)=\int_{\mathbb R}H(t)e^{-2\pi iut}\,dt$
(all formulas follow this normalization).  Xi normalization:
$\xi(s)=\tfrac12s(s-1)\pi^{-s/2}\Gamma(s/2)\zeta(s)$.  Zero counting:
$\rho=\tfrac12+\delta_\rho+i\gamma_\rho$ with multiplicity; sums over all
nontrivial zeros.  Prime-term sign is fixed by the Fourier convention.

\section*{Limitations}
The validity of the criterion depends on the stated analytic framework
and the applicability of the invoked explicit formula.  The equivalence theorem is conditional on the stated distributional
framework and on the validity of the explicitly verified analytic
identities.  Independent verification of the key lemmas by the community
is required before any stronger claim is made.

\end{document}
